\scanpage{38}
\mspage{18}
par une application
\[
(44)\qquad S'\amalg\widetilde A'\longrightarrow S^{\deleted{\unsure{'}}}\amalg\widetilde A^{\deleted{\unsure{'}}}
\]
qui applique $S'$ dans $S^{\deleted{\unsure{'}}}$ (mais pas néc. $\widetilde A'$ dans $\widetilde A^{\deleted{\unsure{'}}}$), qui commute à $\sigma_{G'}$ et $\sigma_{G^{\deleted{\unsure{'}}}}$ (définis via $\sigma_{\widetilde A'}$ et $\sigma_{\widetilde A^{\deleted{\unsure{'}}}}$, avec l'identité sur $S'$ resp. $S^{\deleted{\unsure{'}}}$) et commute à l'application ``origine'' (qui sur $S'$ resp. $S^{\deleted{\unsure{'}}}$ est l'identité).

On peut aussi dire que c'est défini par la donnée
\[
(45)\quad\left\{\begin{aligned}
&\text{1°) d'une application}\qquad f_{\unsure{S|S^0}}:S'\longrightarrow S^{\deleted{\unsure{'}}}\\
&\text{2°) d'une partie }\deleted{\unsure{\widetilde A}}\widetilde B'^{\deleted{\unsure{'}}}\text{ de }\widetilde A',\text{ stable par }\sigma_{\widetilde A'}\\
&\text{3°) d'une application }\widetilde B'\unsure{\to|{-}}\widetilde A^{\deleted{\unsure{'}}}\text{ commutant à }\deleted{\unsure{\sigma}}\\
&\qquad\unsure{\text{à}}\ \sigma_{\widetilde A'}|\widetilde B'\text{ et à }\sigma_{\widetilde A^{\deleted{\unsure{'}}}},\text{ et aux applications}\\
&\qquad o_{\widetilde A'}|\widetilde B'\text{ et }o_{\widetilde A^{\deleted{\unsure{'}}}}.
\end{aligned}\right.
\]
\margin{et il faut supposer de plus que pour $\forall\vec a\in\widetilde A'\setminus\widetilde B'$, l'origine et l'extrémité de $\vec a$ ont $\unsure{\overline m|=}$ image dans $S$ par $f_S$}

On peut dire enfin que c'est aussi la donnée d'un sous-graphe $G_{\widetilde B'}$ de \deleted{\unsure{$G$}}$G'$ ayant $\unsure{=_{\mathrm s}|\overline m_{\mathrm s}}$ sommets (cela \unsure{défini|définit} juste \unsure{rev.|un} à donner une partie $\widetilde B'$ de $\widetilde A'$ stable par $\sigma_{\widetilde A'}$), et un morphisme de $G_{\widetilde B'}$ dans $G^{\deleted{\unsure{'}}}$. La composition des morphismes généralisés est alors définie et \unsure{lue} de façon
\mspage{37}
